\documentclass[a4paper,12pt]{article}
\usepackage[T2A]{fontenc}
\usepackage[utf8]{inputenc}
\usepackage[russian]{babel}
\usepackage{amsmath,amssymb,mathtools}
\renewcommand{\familydefault}{\sfdefault}
\usepackage{icomma}
\usepackage[a4paper,left=23mm,right=23mm,top=22mm,bottom=23mm,headheight=15pt]{geometry}
\usepackage{microtype}
\usepackage{tikz}
\usetikzlibrary{arrows.meta,positioning,shapes.geometric,calc}
\usepackage{tabularx,booktabs}
\usepackage{enumitem}
\usepackage{parskip}
\usepackage{fancyhdr}
\usepackage{setspace}
\usepackage{xcolor}
\usepackage[hidelinks]{hyperref}
\definecolor{accent}{HTML}{087B7D}
\definecolor{darkaccent}{HTML}{125457}
\definecolor{textgray}{HTML}{30363A}
\definecolor{lightline}{HTML}{DAE3E2}
\renewcommand{\today}{23 октября 2026 г.}
\setlength{\parindent}{0pt}
\setlength{\parskip}{6pt}
\onehalfspacing
\setlist[itemize]{leftmargin=*,itemsep=1pt,topsep=2pt}
\setlist[enumerate]{leftmargin=*,itemsep=3pt,topsep=4pt}
\setlength{\emergencystretch}{2em}
\pagestyle{fancy}
\fancyhf{}
\fancyhead[L]{\small\textcolor{darkaccent}{ЧЕК-ЛИСТ · СТЕПЕНИ}}
\fancyhead[R]{\small\textcolor{textgray}{23.10.26}}
\fancyfoot[C]{\small\thepage}
\renewcommand{\headrulewidth}{0.3pt}
\renewcommand{\footrulewidth}{0pt}
\newcommand{\key}[1]{\textcolor{darkaccent}{\textbf{#1}}}
\newcommand{\sub}[1]{\medskip\noindent\textcolor{darkaccent}{\large\bfseries #1}\par\vspace{2pt}}
\tikzset{
 flowstart/.style={ellipse,draw=darkaccent,fill=accent!7,line width=.8pt,minimum height=11mm,minimum width=44mm,align=center,font=\small},
 flowact/.style={rounded corners=2mm,rectangle,draw=darkaccent,fill=white,line width=.7pt,text width=105mm,minimum height=14mm,align=center,inner sep=5pt,font=\small},
 flowdecision/.style={diamond,aspect=3,draw=darkaccent,fill=accent!5,line width=.8pt,text width=44mm,align=center,inner xsep=0pt,inner ysep=2pt,font=\small},
 flowarrow/.style={-{Stealth[length=3mm]},draw=darkaccent,line width=.85pt},
 flowside/.style={rounded corners=2mm,rectangle,draw=darkaccent,fill=white,line width=.7pt,text width=60mm,minimum height=16mm,align=center,inner sep=5pt,font=\small},
 flowchoose/.style={diamond,aspect=2.7,draw=darkaccent,fill=accent!5,line width=.8pt,text width=48mm,align=center,font=\small,inner sep=1pt}
}
\begin{document}
\begin{titlepage}
\thispagestyle{empty}
\begin{center}
\vspace*{30mm}
{\color{darkaccent}\large МАТЕМАТИКА · ПОДГОТОВКА К ПРОФИЛЬНОМУ ЕГЭ}\par
\vspace{22mm}
{\color{darkaccent}\fontsize{34}{39}\selectfont\bfseries Чек-лист}\par
\vspace{10mm}
{\fontsize{22}{28}\selectfont\bfseries Степени и их свойства}\par
\vspace{4mm}
{\large Преобразования · показательные уравнения · дробные показатели}\par
\vspace{23mm}
{\large \today}\par
\vfill
{\normalsize\bfseries Лёвин Артём Александрович}\par
{\normalsize эксклюзивно для Даниила Хрястова}\par
\vspace{19mm}
\begin{tikzpicture}
\draw[accent,line width=1pt] (0,0) .. controls (2,0.7) and (4,-0.65) .. (6,0)
  .. controls (8,0.65) and (10,-0.7) .. (12,0);
\end{tikzpicture}
\vspace{9mm}
\end{center}
\end{titlepage}

\section*{1. Что нужно знать о степенях}
Степень $a^n$ имеет \key{основание} $a$ и \key{показатель} $n$. Для целого положительного $n$ она означает произведение $n$ одинаковых множителей. Если $a\ne0$, то
\[
 a^1=a,\qquad a^0=1,\qquad a^{-n}=\frac{1}{a^n}.
\]
\key{Условие применимости.} Для $a>0$ следующие свойства верны при любых действительных показателях $m,n$. При $a=0$ отрицательные степени не определены. С отрицательным основанием и дробными показателями нужно отдельно проверять существование действительного значения.

\sub{Восемь базовых правил}
\renewcommand{\arraystretch}{1.35}
\begin{tabularx}{\linewidth}{@{}p{6mm}X X@{}}
\toprule
& \textbf{Правило} & \textbf{Как распознать} \\
\midrule
1 & $a^m\cdot a^n=a^{m+n}$ & Основание одинаковое; умножение \\
2 & $\dfrac{a^m}{a^n}=a^{m-n}$ & Основание одинаковое; деление \\
3 & $a^n b^n=(ab)^n$ & Показатель одинаковый; умножение \\
4 & $\dfrac{a^n}{b^n}=\left(\dfrac{a}{b}\right)^n$ & Показатель одинаковый; деление \\
5 & $(a^m)^n=a^{mn}$ & Степень возводится в степень \\
6 & $\dfrac{1}{a^n}=a^{-n}$ & Обратная величина \\
7 & $\left(\dfrac{a}{b}\right)^{-n}=\left(\dfrac{b}{a}\right)^n$ & Смена знака степени и обращение дроби \\
8 & $\sqrt[k]{a^m}=a^{m/k}$ & Корень как дробная степень ($a>0$) \\
\bottomrule
\end{tabularx}
{\small Здесь при наличии деления основания и знаменатели отличны от нуля; в формуле 8 $k\in\mathbb N$, $k\ge2$.}

\key{Читайте равенства в обе стороны.} Например, $6^t7^t=(6\cdot7)^t=42^t$. Общий показатель позволяет объединить основания.

\sub{Сложить или умножить показатели?}
\[
 5^2\cdot5^4=5^{2+4}=5^6,
 \qquad (5^2)^4=5^{2\cdot4}=5^8.
\]
Типичная ошибка --- переносить правило умножения показателей на произведение степеней.

\section*{2. Главный принцип преобразования}
При вычислениях, содержащих степени, ищите возможность перейти к \key{общему основанию}. Удобная последовательность:
\begin{enumerate}
\item \key{Составные числа:} разложите на простые множители и запишите их степени.
\item \key{Корни:} переведите в степени с дробным показателем.
\item \key{Дробные основания:} используйте обратные степени там, где это удобно.
\item \key{Одинаковые основания/показатели:} примените соответствующие правила; проверьте скобки.
\end{enumerate}
Число называют \key{простым}, если оно натуральное и имеет ровно два положительных делителя; \key{составным} --- если их больше двух. Число $1$ не является ни простым, ни составным.
\[
 12=2^2\cdot3,\qquad 72=2^3\cdot3^2,\qquad
 8=2^3,\quad 27=3^3,\quad 25=5^2.
\]
\key{Важно:} составное число не всегда представимо в виде одной степени с целым показателем больше единицы: $12=2^2\cdot3$.

\sub{Небольшая цепочка преобразований}
\[
 \frac{1}{\sqrt[5]{27}}
 =\frac{1}{\sqrt[5]{3^3}}
 =\frac{1}{3^{3/5}}
 =3^{-3/5}.
\]
\[
 0{,}5=\frac12=2^{-1},\qquad
 \sqrt[3]{8}=\sqrt[3]{2^3}=2.
\]

\sub{Скобки: памятка «ВУПС»}
\key{Вычитаемое, Умножаемое, Подставляемое --- Скобки.} Когда на место одного элемента приходит составное выражение, сохраните его границы. При вычитании суммы или разности и при возведении суммы в степень скобки особенно важны:
\[
 5^{\sqrt7-1}\!:\,5^{2-\sqrt7}
 =5^{\,(\sqrt7-1)-(2-\sqrt7)}
 =5^{2\sqrt7-3}.
\]
\[
 \bigl(2^{-1}\bigr)^{2x+1}
 =2^{(-1)(2x+1)}=2^{-2x-1}.
\]
Скобки нужны там, где без них изменяется смысл выражения. При обычном произведении отдельных чисел они могут быть необязательны.

\clearpage
\thispagestyle{plain}
\begin{center}
{\color{darkaccent}\Large\bfseries Алгоритм: приведение к степенному виду}\par
\end{center}
\vspace{7mm}
\begin{center}
\begin{tikzpicture}[node distance=9mm]
\node[flowstart] (s) {Исходное выражение};
\node[flowdecision,below=11mm of s] (d1) {Есть составные\\числа в основаниях?};
\node[flowact,below=13mm of d1] (a1) {Разложите нужные числа на простые множители,\\выделите одинаковые основания};
\node[flowdecision,below=11mm of a1] (d2) {Есть корни\\или дробные основания?};
\node[flowact,below=13mm of d2] (a2) {Запишите корни как дробные степени,\\обратные величины --- как отрицательные степени};
\node[flowact,below=9mm of a2] (a3) {Примените правила степеней;\\сначала проверьте основания, показатели и скобки};
\node[flowstart,below=9mm of a3] (e) {Упрощённое выражение};
\draw[flowarrow] (s)--(d1);
\draw[flowarrow] (d1)--node[right,xshift=2mm]{\small да}(a1);
\draw[flowarrow] (a1)--(d2);
\draw[flowarrow] (d2)--node[right,xshift=2mm]{\small да}(a2);
\draw[flowarrow] (a2)--(a3);
\draw[flowarrow] (a3)--(e);
\draw[flowarrow] (d1.east)--++(30mm,0) |- node[pos=.15,above]{\small нет} (d2.east);
\draw[flowarrow] (d2.west)--++(-30mm,0) |- node[pos=.15,above]{\small нет} (a3.west);
\end{tikzpicture}
\end{center}
\clearpage

\section*{3. Два способа вычислить выражение}
Вычислите
\[
 E=\frac{6^{\sqrt3}\cdot7^{\sqrt3}}{42^{\sqrt3-1}}.
\]
\sub{Способ I. Разложить знаменатель}
\[
\begin{aligned}
 E&=\frac{6^{\sqrt3}\cdot7^{\sqrt3}}{(6\cdot7)^{\sqrt3-1}}
 =\frac{6^{\sqrt3}\cdot7^{\sqrt3}}
 {6^{\sqrt3-1}\cdot7^{\sqrt3-1}}\\
 &=6^{\sqrt3-(\sqrt3-1)}\cdot7^{\sqrt3-(\sqrt3-1)}
 =6\cdot7=42.
\end{aligned}
\]
\sub{Способ II. Свернуть числитель}
\[
 E=\frac{(6\cdot7)^{\sqrt3}}{42^{\sqrt3-1}}
 =42^{\sqrt3-(\sqrt3-1)}=42.
\]
\key{Приём:} ищите общий множитель, общую степень или возможность свести несколько оснований к одному. Выбирайте более короткий путь без нарушения порядка действий.

\section*{4. Показательные уравнения}
\key{Показательное уравнение} --- уравнение, в котором неизвестная входит в показатель степени. Для $a>0$, $a\ne1$ верно:
\[
 \boxed{a^{f(x)}=a^{g(x)}\quad\Longleftrightarrow\quad f(x)=g(x).}
\]
\key{Стратегия:} преобразуйте обе стороны к общему основанию; объедините множители в степень; затем приравняйте показатели.

\sub{Пример 1. Уравнение с множителем}
\[
\begin{aligned}
 2\cdot8^x&=32,\\
 2^1\cdot(2^3)^x&=2^5,\\
 2^{1+3x}&=2^5,\\
 1+3x&=5,\\
 x&=\frac43.
\end{aligned}
\]
Здесь $2^1$ и $2^{3x}$ сначала объединяются. Основание после упрощения должно быть одинаковым в обеих частях.

\sub{Пример 2. Корни в основаниях}
\[
 27\cdot(\sqrt3)^{2x+1}=\frac1{\sqrt[4]{9}}.
\]
\[
\begin{aligned}
 3^3\cdot(3^{1/2})^{2x+1}&=3^{-1/2},\\
 3^{3+\frac12(2x+1)}&=3^{-1/2},\\
 3+x+\frac12&=-\frac12,\\
 x&=-4.
\end{aligned}
\]
Запись $\frac12(2x+1)$ сохраняет обе части показателя: умножается вся скобка.

\section*{5. Уравнения с дробной и отрицательной степенью}
Здесь неизвестная находится \key{в основании}. Важно учитывать, сохраняется ли знак при обратном преобразовании.

\sub{Случай 1. Чётная степень скрыта внутри дробного показателя}
\[
 x^{4/3}=16.
\]
Корень третьей степени определён для всех действительных $x$. Введём $t=\sqrt[3]{x}$. Тогда
\[
 t^4=16\ \Longrightarrow\ t=\pm2
 \Longrightarrow\ \boxed{x=t^3=\pm8}.
\]
\key{Почему корня два?} Формула $(x^{4/3})^{3/4}=|x|$ при действительном $x$ возвращает модуль. Равенство с просто $x$ потеряло бы отрицательное решение.

\sub{Случай 2. Отрицательный чётный показатель}
\[
 x^{-2}=4,\qquad x\ne0.
\]
\[
 \frac1{x^2}=4\ \Longrightarrow\ x^2=\frac14
 \Longrightarrow\ \boxed{x=\pm\frac12}.
\]
Возведение обеих частей в степень $-\frac12$ даёт $|x|$, поэтому знак требуется восстановить.

\sub{Случай 3. Нечётная степень корня}
\[
 x^{5/3}=32\ \Longleftrightarrow\
 (\sqrt[3]{x})^5=32\ \Longleftrightarrow\
 \sqrt[3]{x}=2\ \Longleftrightarrow\ \boxed{x=8}.
\]
Для показателя $5/3$ (после сокращения дроби) функция на всей действительной оси строго возрастает, поэтому здесь один действительный корень.

\textcolor{darkaccent}{\textbf{Самопроверка:}} ОДЗ, знаки, оба корня и подстановка.

\clearpage
\thispagestyle{plain}
\begin{center}
{\color{darkaccent}\Large\bfseries Алгоритм: выбор способа решения уравнения}\par
\end{center}
\vspace{11mm}
\begin{center}
\begin{tikzpicture}[node distance=12mm]
\node[flowstart] (st) {Уравнение со степенями};
\node[flowchoose,below=15mm of st] (q) {Где находится\\неизвестная?};
\node[flowside] (l1) at ($(q.south)+(-42mm,-31mm)$) {\textbf{В показателе}\\Приведите обе части к основанию $a>0$, $a\ne1$};
\node[flowside] (r1) at ($(q.south)+(42mm,-31mm)$) {\textbf{В основании}\\Проверьте допустимость $x$ и вид показателя};
\node[flowside,below=14mm of l1] (l2) {Получите $a^{f(x)}=a^{g(x)}$\\и решите $f(x)=g(x)$};
\node[flowside,below=14mm of r1] (r2) {При чётной степени учтите $\pm$;\\при отрицательной исключите $x=0$};
\node[flowstart] (finish) at ($(l2.south)!0.5!(r2.south)+(0,-22mm)$) {Проверьте найденные корни};
\draw[flowarrow] (st)--(q);
\draw[flowarrow] (q.south west) -| node[pos=.25,above left]{\small показатель}(l1.north);
\draw[flowarrow] (q.south east) -| node[pos=.25,above right]{\small основание}(r1.north);
\draw[flowarrow] (l1)--(l2);
\draw[flowarrow] (r1)--(r2);
\draw[flowarrow] (l2.south) |- (finish.west);
\draw[flowarrow] (r2.south) |- (finish.east);
\end{tikzpicture}
\end{center}
\clearpage

\clearpage
\section*{Тренировочный блок · Путь А}
\textcolor{darkaccent}{\textbf{10 заданий · от базовых преобразований к уравнениям}}\par
Выполняйте задания последовательно. При необходимости пользуйтесь формулами из чек-листа. Решения записывайте с промежуточными преобразованиями.
\begin{enumerate}[label=\textbf{\arabic*.},itemsep=12pt,leftmargin=9mm]
\item Вычислите: \(\displaystyle \frac{7^4\cdot7^2}{7^5}.\)
\item Вычислите: \(\displaystyle \frac{(3^2)^3}{3^4}.\)
\item Разложите число \(108\) на простые множители и запишите результат с использованием степеней.
\item Вычислите: \(\displaystyle \frac{\sqrt[3]{216}}{3^{-1}}.\)
\item Вычислите: \(\displaystyle \frac{(0{,}25)^{-2}}{2^{-1}}.\)
\item Вычислите: \(\displaystyle \frac{4^{\sqrt2}\cdot9^{\sqrt2}}{36^{\sqrt2-1}}.\)
\item Решите уравнение: \(\displaystyle 3\cdot9^x=81.\)
\item Решите уравнение: \(\displaystyle x^{-2}=25.\)
\item Решите уравнение: \(\displaystyle x^{4/3}=256.\)
\item Решите уравнение: \(\displaystyle 125\cdot(\sqrt5)^{2x+3}=\frac1{\sqrt[4]{25}}.\)
\end{enumerate}
\vfill
\noindent\textcolor{darkaccent}{\rule{\linewidth}{0.5pt}}\\[-1mm]
{\small Проверьте ответы только после самостоятельного выполнения всех десяти заданий.}

\clearpage
\section*{Ответы к тренировочному блоку}
\renewcommand{\arraystretch}{1.7}
\begin{center}
\begin{tabularx}{\linewidth}{@{}c X@{}}
\toprule
\textbf{№} & \textbf{Ответ} \\
\midrule
1 & $7$ \\
2 & $9$ \\
3 & $108=2^2\cdot3^3$ \\
4 & $18$ \\
5 & $32$ \\
6 & $36$ \\
7 & $x=\dfrac32$ \\
8 & $x=-\dfrac15$ или $x=\dfrac15$ \\
9 & $x=-64$ или $x=64$ \\
10 & $x=-5$ \\
\bottomrule
\end{tabularx}
\end{center}
\vfill
\begin{center}
\textcolor{darkaccent}{\rule{45mm}{.6pt}}\\
{\small Лёвин Артём Александрович · эксклюзивно для Даниила Хрястова}
\end{center}
\end{document}
